% =============================================================================
% backup.tex — appendix frames (answers to anticipated committee questions)
% Not counted in the 15'. All numbers traceable to results/reports/ & figures/.
% =============================================================================
\begin{frame}{Παράρτημα: διαδικασία τυποποίησης NIST (2016--2025)}
  \centering
  \resizebox{0.9\linewidth}{!}{% VARIANT — not \input by main.tex; promote manually.
% _v1 — beautified original + TOKEN MIGRATION of nist-rounds-flowchart.tex.
% Same vertical top-to-bottom layout & data; local glyph/leader/box styles
% migrated to thesis-tikz-styles.tex tokens. FIX: the original used an UNDEFINED
% `figleader' style for the axis-to-box leaders (no such token in the style
% file) — replaced here with the hairline `figrule, figLW/hair'. The deepening
% process tint (roundbox=#1) is an intentional gradient, kept.
\tikzsetnextfilename{ch1-introduction-nist-rounds-flowchart_v1}
\begin{tikzpicture}[
  vbox/.style={
    draw=figrule, rounded corners=2pt, text width=8.2cm, align=left,
    inner sep=6pt, figT/label, anchor=west
  },
  roundbox/.style={vbox, fill=figprocess!#1},
  milestonebox/.style={vbox, draw=figsolution, figLW/base, fill=figsolutionfill},
  hqcbox/.style={vbox, draw=figaccent, figLW/emph, fill=figaccentfill},
  yr/.style={font=\sffamily\scriptsize\bfseries, text=figneutral, anchor=east},
]
\def\bx{1.3}

% --- Vertical axis ---
\draw[figrule, figLW/emph] (0,0.7) -- (0,-13.0);
\draw[figrule, figLW/emph, -{Stealth[length=5pt]}] (0,-13.0) -- (0,-13.4);

% --- Leader from axis to each box (FIX: was undefined figleader) ---
\foreach \y in {0,-2,-4,-6,-8,-10,-12} { \draw[figrule, figLW/hair] (0,\y) -- (\bx,\y); }

% === Rounds (process, deepening cool tint) ===
\node[figglyph/event] at (0,0) {}; \node[yr] at (-0.18,0) {2016};
\node[roundbox=12, anchor=west] at (\bx,0) {%
  {\sffamily\bfseries\color{figprocess!80!black} Πρόσκληση Υποβολών}\\[1pt]
  Ανοιχτή διαδικασία τυποποίησης PQC· 82 υποβολές $\to$ 69 δεκτοί υποψήφιοι.%
};
\node[figglyph/event] at (0,-2) {}; \node[yr] at (-0.18,-2) {2017};
\node[roundbox=20, anchor=west] at (\bx,-2) {%
  {\sffamily\bfseries\color{figprocess!80!black} Γύρος 1 · Αρχική Αξιολόγηση}\\[1pt]
  Δημόσια κρυπτανάλυση των υποψηφίων· 69 $\to$ 26 για τον Γύρο 2.%
};
\node[figglyph/event] at (0,-4) {}; \node[yr] at (-0.18,-4) {2019};
\node[roundbox=30, anchor=west] at (\bx,-4) {%
  {\sffamily\bfseries\color{figprocess!80!black} Γύρος 2 · Αξιολόγηση Απόδοσης}\\[1pt]
  Απόδοση \& ασφάλεια· 26 $\to$ 7 finalists + 8 alternates.%
};
\node[figglyph/event] at (0,-6) {}; \node[yr] at (-0.18,-6) {2020};
\node[roundbox=40, anchor=west] at (\bx,-6) {%
  {\sffamily\bfseries\color{figprocess!85!black} Γύρος 3 · Τελικοί Υποψήφιοι}\\[1pt]
  7+8 υποψήφιοι $\to$ 4 προς τυποποίηση + 4 KEMs στον Γύρο 4.%
};
\node[figglyph/event] at (0,-8) {}; \node[yr] at (-0.18,-8) {2022};
\node[roundbox=50, anchor=west] at (\bx,-8) {%
  {\sffamily\bfseries\color{figprocess!90!black} Γύρος 4 · Μη-Lattice KEM}\\[1pt]
  Τέσσερις υποψήφιοι: BIKE, Classic McEliece, \HQC, SIKE.%
};

% === Milestones (solution / accent) ===
\node[figglyph/standard] at (0,-10) {}; \node[yr] at (-0.18,-10) {2024};
\node[milestonebox, anchor=west] at (\bx,-10) {%
  {\sffamily\bfseries\color{figsolution!55!black} Πρώτα Πρότυπα FIPS}\\[1pt]
  FIPS 203 (ML-KEM), 204 (ML-DSA), 205 (SLH-DSA) — lattice \& hash-based.%
};
\node[figglyph/focal] at (0,-12) {}; \node[yr] at (-0.18,-12) {2025};
\node[hqcbox, anchor=west] at (\bx,-12) {%
  {\sffamily\bfseries\color{figaccent} Επιλογή \HQC · Μάρτιος 2025}\\[1pt]
  Πέμπτο PQC πρότυπο· draft FIPS 2026, τελικό πρότυπο 2027.%
};

\end{tikzpicture}
}
  % src: figures/ch1-introduction/nist-rounds-flowchart.tex
\end{frame}

\begin{frame}{Παράρτημα: παράμετροι \HQC}
  \centering
  \renewcommand{\arraystretch}{1.2}
  {\small
  \begin{tabular}{@{}lrrrr@{}}
    \toprule
    Επίπεδο & $N$ & $n_1$ & $n_2$ & $\omega$ \\
    \midrule
    \HQC-128 & \num{17669} & 46 & 384 & 66 \\
    \HQC-192 & \num{35851} & 56 & 640 & 100 \\
    \HQC-256 & \num{57637} & 90 & 640 & 131 \\
    \bottomrule
  \end{tabular}}
  % src: figures/ch2-background/hqc-params-table.tex
  \par\medskip
  \begin{minipage}{0.86\textwidth}
    \begin{itemize}
      \item $N$ : διάσταση δακτυλίου (μέγιστος βαθμός πολυωνύμου)
      \item $n_1$ : μήκος \en{RS} codeword (σύμβολα)
      \item $n_2$ : μήκος \en{RM} codeword (bits)
      \item $\omega$ : βάρος \en{Hamming} μυστικού
    \end{itemize}
  \end{minipage}
\end{frame}

\begin{frame}{Παράρτημα: \en{Keccak} (sponge \& state) και οι πράξεις \Zbb}
  \begin{columns}[T,onlytextwidth]
    \begin{column}{0.50\textwidth}
      \resizebox{\linewidth}{!}{% Sponge construction diagram for SHA-3/Keccak
% Shows absorb phase (XOR + permute) and squeeze phase (extract + permute)
\begin{tikzpicture}[
  >=stealth,
  block/.style={draw, minimum width=1.8cm, minimum height=0.9cm,
                font=\small\sffamily, align=center, rounded corners=2pt},
  xornode/.style={draw, circle, inner sep=1.5pt, font=\scriptsize},
  arr/.style={->, thick},
  brace/.style={decorate, decoration={brace, amplitude=6pt, raise=2pt}},
  label/.style={font=\footnotesize\sffamily, text=gray!70!black}
]

% State representation
\def\stateW{3.0}
\def\rateH{1.2}
\def\capH{0.6}

% --- Absorb phase ---
\node[label] at (0, 2.5) {\textbf{Absorb}};

% Input blocks
\node[font=\small] (m0) at (-1.5, 1.5) {$m_0$};
\node[font=\small] (m1) at (2.0, 1.5) {$m_1$};
\node[font=\small] (m2) at (5.5, 1.5) {$m_2$};
\node[font=\small] (dots) at (8.2, 1.5) {$\cdots$};

% XOR nodes
\node[xornode] (xor0) at (-1.5, 0.5) {$\oplus$};
\node[xornode] (xor1) at (2.0, 0.5) {$\oplus$};
\node[xornode] (xor2) at (5.5, 0.5) {$\oplus$};

% Permutation blocks
\node[block, fill=blue!10] (f0) at (0.2, -0.5) {$f$};
\node[block, fill=blue!10] (f1) at (3.7, -0.5) {$f$};
\node[block, fill=blue!10] (f2) at (7.2, -0.5) {$f$};

% State boxes (rate + capacity)
\draw[thick, fill=green!8] (-2.5, 0.0) rectangle (-2.5+\stateW, \rateH);
\node[font=\scriptsize] at (-2.5+\stateW/2, \rateH/2) {$r$ bits};
\draw[thick, fill=red!8] (-2.5, -\capH) rectangle (-2.5+\stateW, 0.0);
\node[font=\scriptsize] at (-2.5+\stateW/2, -\capH/2) {$c$ bits};

% Arrows: input to XOR
\draw[arr] (m0) -- (xor0);
\draw[arr] (m1) -- (xor1);
\draw[arr] (m2) -- (xor2);

% Arrows: XOR to f
\draw[arr] (xor0) -- (f0);
\draw[arr] (xor1) -- (f1);
\draw[arr] (xor2) -- (f2);

% Arrows: f to next XOR
\draw[arr] (f0) -| (xor1);
\draw[arr] (f1) -| (xor2);
\draw[arr] (f2) -- (8.2, -0.5);

% Initial state arrow
\draw[arr] (-2.5+\stateW, 0.3) -- (xor0);

% --- Squeeze phase ---
\node[label] at (9.5, 2.5) {\textbf{Squeeze}};

\node[block, fill=blue!10] (fs0) at (10.0, -0.5) {$f$};
\node[block, fill=blue!10] (fs1) at (13.0, -0.5) {$f$};

% Output blocks
\node[font=\small] (z0) at (10.0, 1.5) {$z_0$};
\node[font=\small] (z1) at (13.0, 1.5) {$z_1$};
\node[font=\small] (zdots) at (14.8, 1.5) {$\cdots$};

% Arrows: f to output
\draw[arr] (fs0) -- (z0);
\draw[arr] (fs1) -- (z1);

% Arrows: chain
\draw[arr] (8.2, -0.5) -- (fs0);
\draw[arr] (fs0) -- (fs1);
\draw[arr] (fs1) -- (14.8, -0.5);

% Rate/capacity annotation
\draw[brace] (-2.5, \rateH+0.1) -- node[above=8pt, font=\scriptsize\sffamily]
  {rate $r$} (-2.5+\stateW, \rateH+0.1);
\draw[brace] (-2.5+\stateW, -\capH-0.1) -- node[below=8pt, font=\scriptsize\sffamily]
  {capacity $c$} (-2.5, -\capH-0.1);

% Annotation: state width
\node[font=\scriptsize\sffamily, text=gray] at (-1.0, -1.5)
  {state = $r + c = \num{1600}$ bits};

\end{tikzpicture}
}
    \end{column}
    \begin{column}{0.04\textwidth}
      \centering
      {\tikzexternaldisable
       \begin{tikzpicture}
         \draw[dashed, gray!60, line width=0.6pt] (0,0) -- (0,3.8);
       \end{tikzpicture}}
    \end{column}
    \begin{column}{0.42\textwidth}
      \resizebox{\linewidth}{!}{\tikzsetnextfilename{ch2-background-keccak-state-3d}
\begin{tikzpicture}[
  x={(0.9cm, 0cm)}, y={(0.35cm, 0.35cm)}, z={(0cm, 0.8cm)},
  font=\sffamily\small
]

\def\cellsize{0.8}
\def\depth{2.5}

% Bottom face grid
\foreach \x in {0,...,4} {
  \foreach \y in {0,...,4} {
    \draw[figLW/hair, draw=figrule, fill=black!4, fill opacity=0.5]
      (\x*\cellsize, \y*\cellsize, 0) --
      (\x*\cellsize+\cellsize, \y*\cellsize, 0) --
      (\x*\cellsize+\cellsize, \y*\cellsize+\cellsize, 0) --
      (\x*\cellsize, \y*\cellsize+\cellsize, 0) -- cycle;
  }
}

% Lane highlight (x=1, y=2) — figaccent (three visible faces)
\fill[figaccent, opacity=0.35]
  (1*\cellsize, 2*\cellsize, 0) -- (1*\cellsize+\cellsize, 2*\cellsize, 0) --
  (1*\cellsize+\cellsize, 2*\cellsize, \depth) -- (1*\cellsize, 2*\cellsize, \depth) -- cycle;
\fill[figaccent, opacity=0.25]
  (1*\cellsize, 2*\cellsize, 0) -- (1*\cellsize, 2*\cellsize+\cellsize, 0) --
  (1*\cellsize, 2*\cellsize+\cellsize, \depth) -- (1*\cellsize, 2*\cellsize, \depth) -- cycle;
\fill[figaccent, opacity=0.45]
  (1*\cellsize, 2*\cellsize, \depth) -- (1*\cellsize+\cellsize, 2*\cellsize, \depth) --
  (1*\cellsize+\cellsize, 2*\cellsize+\cellsize, \depth) -- (1*\cellsize, 2*\cellsize+\cellsize, \depth) -- cycle;

% Column highlight (x=3, all y) — figsolution
\foreach \y in {0,...,4} {
  \fill[figsolution, opacity=0.30]
    (3*\cellsize, \y*\cellsize, 0) -- (3*\cellsize+\cellsize, \y*\cellsize, 0) --
    (3*\cellsize+\cellsize, \y*\cellsize+\cellsize, 0) -- (3*\cellsize, \y*\cellsize+\cellsize, 0) -- cycle;
}

% Vertical depth edges
\foreach \x in {0,4} {
  \foreach \y in {0,4} {
    \draw[figLW/hair, draw=figrule] (\x*\cellsize, \y*\cellsize, 0) -- (\x*\cellsize, \y*\cellsize, \depth);
  }
}
\draw[figLW/hair, draw=figrule] (5*\cellsize, 0, 0) -- (5*\cellsize, 0, \depth);
\draw[figLW/hair, draw=figrule] (5*\cellsize, 5*\cellsize, 0) -- (5*\cellsize, 5*\cellsize, \depth);
\draw[figLW/hair, draw=figrule] (0, 5*\cellsize, 0) -- (0, 5*\cellsize, \depth);

% Top face outline
\draw[figLW/hair, draw=figrule]
  (0, 0, \depth) -- (5*\cellsize, 0, \depth) --
  (5*\cellsize, 5*\cellsize, \depth) -- (0, 5*\cellsize, \depth) -- cycle;

% Axis labels
\draw[figLW/base, draw=figneutral, <->] (-0.5, 0, 0) -- (-0.5, 5*\cellsize, 0)
  node[midway, left, figT/micro] {$y$ (5)};
\draw[figLW/base, draw=figneutral, <->] (0, -0.5, 0) -- (5*\cellsize, -0.5, 0)
  node[midway, below, figT/micro] {$x$ (5)};
\draw[figLW/base, draw=figneutral, <->] (5*\cellsize+0.5, 0, 0) -- (5*\cellsize+0.5, 0, \depth)
  node[midway, right, figT/micro] {$z$ (64 bits)};

% Annotations
\node[figT/micro, text=figaccent!70!black, anchor=west] at (6, 2.5*\cellsize, 1.2)
  {Lane $A[x,y]$};
\draw[figLW/hair, draw=figaccent!70!black, ->] (6, 2.5*\cellsize, 1.2) -- (1.5*\cellsize, 2.5*\cellsize, 1.2);

\node[figT/micro, text=figsolution!60!black, anchor=west] at (6, 0.5*\cellsize, 0)
  {Column $A[x,*,z]$};
\draw[figLW/hair, draw=figsolution!60!black, ->] (6, 0.5*\cellsize, 0) -- (3.5*\cellsize, 2*\cellsize, 0);

% Title annotation — dropped straight down (negative z) for clear separation from the figure
\node[figT/annot, anchor=north] at (2.8*\cellsize, 0, -1.0)
  {Κατάσταση Keccak: $5 \times 5 \times 64 = \num{1600}$ bits};

\end{tikzpicture}
}
    \end{column}
  \end{columns}
  \vspace{0.3em}
  \begin{itemize}
    \item \defterm{absorb}: \en{XOR} εισόδου στο rate, μετάθεση $f$
    \item \defterm{squeeze}: εξαγωγή εξόδου από την ίδια κατάσταση
    \item bit-level: rotation lane, εναλλαγή lanes, αντιστοίχιση σε \Zbb (\rori, \en{andn})
  \end{itemize}
\end{frame}

\begin{frame}{Παράρτημα: συνενωμένη κωδικοποίηση (\en{RS} $\circ$ \en{RM})}
  \centering
  \resizebox{0.75\linewidth}{!}{% Concatenated ECC message journey: encode (RS -> RM) -> noisy channel -> decode (RM -> RS)
% Serpentine layout: top row left-to-right (encode), bottom row right-to-left (decode).
\tikzsetnextfilename{ch2-background-concatenated-pipeline}
\begin{tikzpicture}[
  >=stealth,
  block/.style={rectangle, draw, fill=#1, minimum width=2.7cm,
    minimum height=1.15cm, text width=2.5cm, align=center,
    rounded corners=2pt, font=\small},
  block/.default={blue!8},
  arrow/.style={->, thick},
  lbl/.style={font=\scriptsize, midway, above},
  vlbl/.style={font=\scriptsize, midway, right},
  ann/.style={font=\scriptsize, text=gray},
  phase/.style={font=\sffamily\bfseries, text=thesisblue},
]

% Phase headers
\node[phase] at (6, 1.5) {Κωδικοποίηση (\Encaps)};
\node[phase] at (6, -4.7) {Αποκωδικοποίηση (\Decaps)};

% Row 1 -- encode (left -> right)
\node[block=blue!10]   (m0)   at (0,0)  {Μήνυμα $m$\\($k$ bytes)};
\node[block=green!12]  (rse)  at (4,0)  {Reed-Solomon Κωδικοποιητής (εξωτερικός)};
\node[block=orange!15] (rme)  at (8,0)  {RM(1,7) Κωδικοποιητής (εσωτερικός)};
\node[block=red!12]    (chan) at (12,0) {Κανάλι / θόρυβος \Decaps};

% Row 2 -- decode (right -> left)
\node[block=yellow!15] (recv) at (12,-3) {Ληφθείσα λέξη ($n_1 \times n_2$ bits)};
\node[block=orange!15] (rmd)  at (8,-3)  {RM(1,7) Αποκωδικοποιητής (εσωτερικός)};
\node[block=green!12]  (rsd)  at (4,-3)  {Reed-Solomon Αποκωδικοποιητής (εξωτερικός)};
\node[block=blue!10]   (mhat) at (0,-3)  {Ανακτημένο μήνυμα $\hat{m}$ ($k$ bytes)};

% Encode arrows (left -> right)
\draw[arrow] (m0)  -- (rse)  node[lbl] {$k$ bytes};
\draw[arrow] (rse) -- (rme)  node[lbl] {$n_1$ σύμβολα};
\draw[arrow] (rme) -- (chan) node[lbl] {$n_1 \times n_2$ bits};

% Serpentine link down
\draw[arrow] (chan) -- (recv) node[vlbl] {\Decaps};

% Decode arrows (right -> left)
\draw[arrow] (recv) -- (rmd) node[lbl] {$n_1 \times n_2$ bits};
\draw[arrow] (rmd)  -- (rsd) node[lbl] {$n_1$ σύμβολα};
\draw[arrow] (rsd)  -- (mhat) node[lbl] {$k$ bytes};

% Granularity annotations (under decode row)
\node[ann, below=0.25cm of rmd] {επίπεδο bit};
\node[ann, below=0.25cm of rsd] {επίπεδο συμβόλου};

\end{tikzpicture}
}
  \begin{itemize}
    \item εσωτερικός \en{Reed--Muller}: διορθώνει σφάλματα bit εντός κάθε συμβόλου
    \item εξωτερικός \en{Reed--Solomon}: διορθώνει ολόκληρα σύμβολα (bytes) που ο εσωτερικός απέτυχε να αποκαταστήσει
  \end{itemize}
\end{frame}

\begin{frame}{Παράρτημα: επεκτάσεις RISC-V και προφίλ \en{RVA23}}
  \centering
  % figures/ch3-riscv/rva23-extensions-tables.tex
% Class-agnostic shared tables (macros only, no \gls/\cite) so both the thesis
% and the defense deck can \input the same source. Values traceable to:
% src: chapters/03-risc-v-architecture.tex (tab:extensions-used, tab:rva23-extensions)
\begingroup
\small
\begin{tblr}{
  width = \linewidth,
  colspec = {l X[l] l c},
  row{1} = {font=\bfseries},
  hline{1,Z} = {0.6pt},
  hline{2} = {0.4pt},
}
  Επέκταση & Περιγραφή & Ratified & Tier \\
  \Zbc  & Carry-less multiply (\clmul, \clmulh) & 11/2021 & 1 \\
  \Zbb  & Basic bit-manip (\rori, \insn{andn}, \insn{rev8}) & 11/2021 & 1 \\
  V     & Vector extension 1.0 (\RVV) & 11/2021 & 2 \\
  \Zvbc & Vector carry-less multiply (\vclmul, \vclmulh) & 09/2023 & 3 \\
  \Zvbb & Vector bit-manip (\vror, \insn{vandn}, \insn{vbrev8}) & 09/2023 & 3 \\
\end{tblr}

\vspace{1em}

\begin{tblr}{
  width = \linewidth,
  colspec = {l l X[l]},
  row{1} = {font=\bfseries},
  hline{1,Z} = {0.6pt},
  hline{2} = {0.4pt},
}
  Επέκταση & Κατάσταση RVA23 & Σημασία \\
  V (\RVV\ 1.0) & Mandatory & Εγγυημένη παρουσία σε κάθε RVA23 \en{chip} \\
  \Zbb  & Mandatory & Rotations, \insn{andn} διαθέσιμα παντού \\
  \Zvbb & Mandatory & Vector rotations/\insn{vandn} εγγυημένα \\
  \Zvbc & Development option & Αναμένεται υποχρεωτική σε μελλοντικό προφίλ \\
  \Zbc  & Expansion option & Scalar carry-less multiply \\
\end{tblr}
\endgroup

  % src: chapters/03-risc-v-architecture.tex (tab:extensions-used, tab:rva23-extensions)
\end{frame}

\begin{frame}{Παράρτημα: \en{RVV} και οι παράμετροι \VLEN/\SEW/\LMUL/\VLMAX}
  \begin{columns}[T,onlytextwidth]
    \begin{column}{0.55\textwidth}
      \resizebox{\linewidth}{!}{\tikzsetnextfilename{ch3-riscv-rvv-lmul-grouping_emphasis}
\begin{tikzpicture}[
  reg/.style={draw=figrule, minimum width=1.8cm, minimum height=0.32cm,
              font=\ttfamily\tiny},
  annot/.style={figT/annot},
  grouphdr/.style={figT/title},
  group/.style={rounded corners=2pt, inner sep=2pt}
]

% === LMUL = 1 (left) ===
\node[grouphdr] at (0, 1.0) {LMUL = 1};
\node[annot] at (0, 0.6) {32 ανεξάρτητοι καταχωρητές};
\foreach \i in {0,...,7} {
  \pgfmathsetmacro{\ypos}{-\i * 0.4}
  \node[reg, fill=figprocess!10] (l1v\i) at (0, \ypos) {v\i};
}
\node[annot] at (0, -3.5) {$\vdots$};
\node[annot] at (0, -4.0) {VLMAX = VLEN/SEW};

% === LMUL = 4 (right) — emphasised groups ===
\node[grouphdr] at (5, 1.0) {LMUL = 4};
\node[annot] at (5, 0.6) {8 ομάδες × 4 καταχωρητές};

% Group 0: v0-v3
\foreach [count=\idx from 0] \i in {0,1,2,3} {
  \pgfmathsetmacro{\ypos}{-\idx * 0.4}
  \node[reg, fill=figprocess!15] (l4v\i) at (5, \ypos) {v\i};
}
\node[group, draw=figprocess, figLW/emph, fit=(l4v0)(l4v3)] (grp0) {};
\node[figpill, anchor=west] at (6.35, -0.6) {group 0};

% Group 1: v4-v7
\foreach [count=\idx from 0] \i in {4,5,6,7} {
  \pgfmathsetmacro{\ypos}{-1.9 - \idx * 0.4}
  \node[reg, fill=figsolution!15] (l4v\i) at (5, \ypos) {v\i};
}
\node[group, draw=figsolution, figLW/emph, fit=(l4v4)(l4v7)] (grp1) {};
\node[figpill, anchor=west] at (6.35, -2.5) {group 1};

% Brace: one logical operand spans 4 registers
\draw[figbrace] (grp0.north east) -- (grp0.south east);
% (1) operand explainer moved onto group 0's row, right of the pill
\node[annot, anchor=west, text width=3.0cm] at (7.95, -0.6)
  {1 logical operand =\\4 registers};

\node[annot] at (5, -3.5) {$\vdots$};
\node[annot] at (5, -4.0) {VLMAX = (VLEN/SEW) $\times$ LMUL};

% (2) Trade-off annotation — widened + no hyphenation -> single clean line
\node[annot, text width=13cm, align=left, anchor=north west] at (-1.5, -4.8)
  {\hyphenpenalty=10000\exhyphenpenalty=10000
   Αύξηση LMUL: περισσότερα στοιχεία ανά εντολή $\leftrightarrow$ λιγότεροι διαθέσιμοι ανεξάρτητοι τελεστέοι (register pressure).};

\end{tikzpicture}
}
    \end{column}
    \begin{column}{0.45\textwidth}
      \footnotesize
      \begin{itemize}
        \item \HL{\VLEN}: πλάτος φυσικού \en{vector register} σε bits, \en{hardware-defined} (π.χ. 128/256/512), σταθερό στο runtime
        \item \HL{\SEW}: πλάτος ενός στοιχείου σε bits (8/16/32/64), ρυθμιζόμενο runtime μέσω \insn{vsetvli}
        \item \HL{\LMUL}: πόσα registers ομαδοποιούνται σε ένα λογικό \en{vector group} (1, 2, 4, 8 ή κλασματικά 1/2, 1/4, 1/8)
        \item \HL{\VLMAX}: μέγιστος αριθμός στοιχείων ανά \en{vector operation} στα τρέχοντα \SEW/\LMUL
      \end{itemize}
      \vspace{0.4em}
      \[ \VLMAX = \frac{\VLEN \times \LMUL}{\SEW} = \frac{256 \times 4}{64} = 16 \]
      {\scriptsize \VLEN=256, \SEW=64, \LMUL=4: μία εντολή επεξεργάζεται έως 16 στοιχεία των 64 bits.}
    \end{column}
  \end{columns}
\end{frame}

\begin{frame}{Παράρτημα: το ερευνητικό κενό}
  \centering
  \adjustbox{max width=\linewidth,max totalheight=0.72\textheight,center}{% Slide variant (trimmed) of \label{tab:rw-riscv-summary}.
% Full 7-col table: chapters/04-related-work.tex (canonical source).
% Dropped cols "Μέθοδος polymul" + "Στρατηγική επιτάχυνσης" -> frame \pdfpcnote.
\renewcommand{\arraystretch}{1.15}%
\begin{tabular}{@{}ll ccc@{}}
  \toprule
  Εργασία & Platform & \rotatebox{70}{RISC-V} & \rotatebox{70}{Καθαρό λογισμ.} & \rotatebox{70}{Πλήρες \KEM} \\
  \midrule
  \multicolumn{5}{@{}l}{\textit{\HQC σε RISC-V (custom hardware), το ερευνητικό κενό}} \\
  \en{Sch\"offel} et al. & RV32IC & \ding{51} & \ding{55} & \ding{51} \\
  \en{Ras} et al. & RV32IMC & \ding{51} & \ding{55} & \ding{51} \\
  \en{TYRCA} & CV32E40PX & \ding{51} & \ding{55} & \ding{51} \\
  \en{Seo} et al. & RV32 (Ibex) & \ding{51} & \ding{55} & \ding{51} \\
  \midrule
  \multicolumn{5}{@{}l}{\textit{Υλοποιήσεις καθαρού λογισμικού σε x86 / ARM}} \\
  \en{Robert \& V\'eron} & x86 (AVX-512/AVX2) & \ding{55} & \ding{51} & \ding{55} \\
  \en{Jang} et al. & x86 / ARM & \ding{55} & \ding{51} & \ding{55} \\
  \en{Chen} et al. & x86 / ARM & \ding{55} & \ding{51} & \ding{51} \\
  \en{OptHQC} & x86 AVX2 & \ding{55} & \ding{51} & \ding{51} \\
  \midrule
  \multicolumn{5}{@{}l}{\textit{Λοιπές εργασίες (ISE / lattice)}} \\
  \en{HORCRUX} & RISC-V + ISE & \ding{51} & \ding{55} & — \\
  \en{ECO-CRYSTALS} & RV64IMB & \ding{51} & \ding{51} & — \\
  \en{Zhang} et al. & RV\{32,64\}IM\{B\}\{V\} & \ding{51} & \ding{51} & — \\
  \midrule
  \HL{Παρούσα εργασία} & \HL{RV64GC\{V\}} & \ding{51} & \ding{51} & \ding{51} \\
  \bottomrule
\end{tabular}%
}
  % src: chapters/04-related-work.tex :: tab:rw-riscv-summary (full 7-col table)
  \pdfpcnote{gap: όλες οι RISC-V εργασίες = custom hardware (loosely/tightly-coupled ISA) ή δεν καλύπτουν πλήρες KEM. Παρούσα εργασία: επικυρωμένες επεκτάσεις (Zbc/Zvbc), πλήρες KEM.}
  \pdfpcnote{}
  \pdfpcnote{method/strategy ammo (dropped columns):}
  \pdfpcnote{- Schoffel: sparse x dense, loosely-coupled + custom ISA}
  \pdfpcnote{- Ras (ANDROMEDA): FAFFT, loosely-coupled}
  \pdfpcnote{- TYRCA: Karatsuba, tightly-coupled (CV-X-IF)}
  \pdfpcnote{- Seo (Ibex): Toom-Cook + Karatsuba, tightly-coupled}
  \pdfpcnote{- Robert/Veron, Jang: Toom-Cook + Karatsuba, SIMD (x86/ARM)}
  \pdfpcnote{- Chen: additive FFT (CRT), SIMD. OptHQC: sparse x dense, AVX2}
  \pdfpcnote{- HORCRUX/ECO-CRYSTALS/Zhang: NTT / custom PQC ISE (lattice/ISE), oxi HQC full KEM}
  \pdfpcnote{- Parousa ergasia: Karatsuba, epikyromeno SW (Zbc/Zvbc)}
\end{frame}

\begin{frame}{Παράρτημα: γιατί ISA Extensions και C Intrinsics}
  \pdfpcnote{budget on-demand | trigger: "γιατί ISA/intrinsics, όχι FPGA/ASIC ή assembly"}
  \pdfpcnote{}
  \pdfpcnote{- άξονας υλικού → ISA ext = sweet spot → κέρδος χωρίς κόστος/χρόνο σχεδίασης FPGA/ASIC}
  \pdfpcnote{- άξονας λογισμικού → intrinsics → έλεγχος register/instruction + βελτιστοποιήσεις compiler (scheduling, regalloc)}
  \pdfpcnote{- κοινό μήνυμα → σκόπιμο ενδιάμεσο, portable + maintainable}
  \pdfpcnote{}
  \pdfpcnote{cut: δείξε μόνο τα δύο pills («Παρούσα εργασία») και προχώρα}
  \centering
  \resizebox{0.82\textwidth}{!}{% VARIANT — not \input by main.tex; promote manually.
% _v1 — beautified original + TOKEN MIGRATION of acceleration-spectrum.tex.
% Same 4-box flexibility<->performance layout; inline fonts/weights migrated to
% thesis-tikz-styles.tex tokens, the "you are here" pill uses the figpill token,
% and the focal box uses figbox/focal weights. FIX: the original used an
% UNDEFINED `figaccentleader' style for the pill connector (no such token) —
% replaced with `draw=figaccent, figLW/base'. Kept inside \resizebox so the wide
% layout never overflows the text block.
\resizebox{\textwidth}{!}{%
\tikzsetnextfilename{ch1-introduction-acceleration-spectrum_v1}
\begin{tikzpicture}[
  node distance=0.8cm,
  box/.style={
    draw=figrule, rounded corners=3pt, minimum width=3cm, minimum height=1.4cm,
    align=center, figT/label, fill=white
  },
  highlight/.style={box, draw=figaccent, figLW/emph, fill=figaccentfill},
  endlabel/.style={figT/annot},
  rowtitle/.style={font=\sffamily\footnotesize\bfseries, text=figneutral, anchor=west},
  ch/.style={figT/annot, align=center},
]

% --- Boxes (left = most flexible, right = least flexible) ---
\node[box] (sw) {Γενικός κώδικας C\\[-2pt]\footnotesize (base ISA)};
\node[highlight, right=0.8cm of sw] (isa) {ISA Extensions\\[-2pt]\footnotesize (Zbc, Zvbc, Zvbb)};
\node[box, right=0.8cm of isa] (fpga) {FPGA\\[-2pt]\footnotesize Accelerator};
\node[box, right=0.8cm of fpga] (asic) {Custom ASIC\\[-2pt]\footnotesize (σταθερή λειτουργία)};

% --- Gradient spectrum bar (flexibility[green] -> performance[vermillion]) ---
\shade[left color=figsolution!35, right color=figthreat!42, rounded corners=2pt]
  ([yshift=0.70cm]sw.north west) rectangle ([yshift=1.18cm]asic.north east);
\node[endlabel, anchor=west] at ([xshift=2pt,yshift=0.94cm]sw.north west)
  {$\leftarrow$ Μεγαλύτερη ευελιξία};
\node[endlabel, anchor=east] at ([xshift=-2pt,yshift=0.94cm]asic.north east)
  {Μεγαλύτερη απόδοση $\rightarrow$};

% --- "You are here" pill on the ISA box (figpill token) ---
\node[figpill, below=0.55cm of isa] (here) {Παρούσα εργασία};
\draw[draw=figaccent, figLW/base] (isa.south) -- (here.north);

% --- "Χρόνος αλλαγής": row heading + per-box values ---
\coordinate (titley) at ([yshift=-1.9cm]sw.south);
\node[rowtitle] at (sw.west |- titley) {Χρόνος αλλαγής:};
\coordinate (rowy) at ([yshift=-2.5cm]sw.south);
\node[ch] at (sw |- rowy)   {λεπτά–ώρες};
\node[ch] at (isa |- rowy)  {λεπτά (recompile)};
\node[ch] at (fpga |- rowy) {ημέρες–εβδομάδες};
\node[ch] at (asic |- rowy) {μήνες–έτη};

\end{tikzpicture}%
}% end resizebox
}
  % src: figures/ch1-introduction/acceleration-spectrum.tex
  \par\medskip
  \textcolor{gray!45}{\rule{0.55\linewidth}{0.4pt}}
  \par\medskip
  \resizebox{0.82\textwidth}{!}{% figures/riscv-programming-spectrum_gradient.tex
% VARIANT (_gradient, ch1 house style) — not \input by main.tex; promote manually.
% Mirrors docs/thesis/figures/ch1-introduction/acceleration-spectrum.tex:
% white boxes chained left->right, a single shaded gradient bar on top encoding
% the control<->optimisation trade-off, and a «Παρούσα εργασία» figpill on the
% method this thesis actually uses (C Intrinsics for the crypto kernels).
% NOTE: move the pill to another box if the «current work» method differs.
\resizebox{\textwidth}{!}{%
\tikzsetnextfilename{ch3-riscv-riscv-programming-spectrum_gradient}
\begin{tikzpicture}[
    node distance=0.8cm,
    box/.style={
      draw=figrule, rounded corners=3pt, minimum width=3cm, minimum height=1.4cm,
      align=center, figT/label, fill=white
    },
    highlight/.style={box, draw=figaccent, figLW/emph, fill=figaccentfill},
    endlabel/.style={figT/annot},
  ]

  % Boxes (left = most control, right = most compiler optimisation)
  \node[box]                              (asm)    {Pure Assembly\\[-2pt]\footnotesize (\texttt{.S})};
  \node[box, right=0.8cm of asm]          (inline) {Inline Assembly\\[-2pt]\footnotesize (\texttt{\_\_asm\_\_})};
  \node[highlight, right=0.8cm of inline] (intr)   {C Intrinsics\\[-2pt]\footnotesize (\texttt{\_\_riscv\_*})};
  \node[box, right=0.8cm of intr]         (auto)   {Auto-Lowering\\[-2pt]\footnotesize (Plain C)};

  % Gradient spectrum bar: control[vermillion] -> optimisation[green]
  \shade[left color=figthreat!42, right color=figsolution!35, rounded corners=2pt]
    ([yshift=0.70cm]asm.north west) rectangle ([yshift=1.18cm]auto.north east);
  \node[endlabel, anchor=west] at ([xshift=2pt,yshift=0.94cm]asm.north west)
    {$\leftarrow$ Μεγαλύτερος έλεγχος};
  \node[endlabel, anchor=east] at ([xshift=-2pt,yshift=0.94cm]auto.north east)
    {Βελτιστοποίηση compiler $\rightarrow$};

  % «You are here» pill on the C Intrinsics box
  \node[figpill, below=0.55cm of intr] (here) {Παρούσα εργασία};
  \draw[draw=figaccent, figLW/base] (intr.south) -- (here.north);

\end{tikzpicture}%
}% end resizebox
}
  % src: figures/ch3-riscv/riscv-programming-spectrum.tex
\end{frame}

\begin{frame}[fragile]{Παράρτημα: \en{Zbb} idiom recognition στον \en{Keccak}}
  Ο μεταγλωττιστής (\en{Clang}~21.1.1, \code{-O3}) αναγνωρίζει αντιπροσωπευτικά
  \en{C} idioms της μετάθεσης και εκπέμπει μονολεκτικές εντολές \Zbb:
  \begin{lstlisting}
(x << n) | (x >> (64-n))  =>  rori   // rotation idiom, order-independent
(~a) & b                  =>  andn   // not+and fused into one op
  \end{lstlisting}
  \begin{itemize}\small
    \item \pos{Αποτέλεσμα}: ${\approx}$\speedup{1.36} στον πυρήνα \en{Keccak}-$f[1600]$ (μόνο \code{-march}, χωρίς αλλαγή πηγαίου)
  \end{itemize}
\end{frame}

\begin{frame}{Παράρτημα: διανυσματοποίηση \en{decode} (\en{RM} FHT $+$ \en{RS})}
  \begin{itemize}\small
    \item \code{expand\_and\_sum} (bit-extraction): \insn{vid} για shift amounts $[0,\ldots,31]$, \insn{vsrl.vv}$+$\insn{vand} εξάγει \num{32} bits/word (\en{SEW=32}, \en{LMUL=4})
    \item \code{hadamard} (7-pass FHT): \num{64} butterfly ζεύγη με \insn{vlse16} (strided, byte-stride \num{4}), \num{6} vector εντολές/pass έναντι \num{896} scalar 
    \item \code{find\_peaks} (argmax): \insn{vredmaxu}$+$\insn{vmseq}/\insn{vfirst\_m} σε ένα pass (\en{LMUL=8})
    \item \en{RS} Στάδιο~A (syndromes): dot-product επί του \GF{2^8}
  \end{itemize}
\end{frame}

\begin{frame}{Παράρτημα: βαθμωτή κλιμάκωση ανά στάδιο (\en{ablation}, \instret)}
  \small
  \begin{center}
  \begin{tabular}{@{}lccc@{}}
    \toprule
    Στάδιο (\instret) & \HQC-128 & \HQC-192 & \HQC-256 \\
    \midrule
    A: word-level \clmul & \speedup{38.6} & \speedup{38.4} & \speedup{45.4} \\
    B: static tree & \speedup{1.99} & \speedup{1.99} & \speedup{2.05} \\
    C: assign-on-first-write & \speedup{1.27} & \speedup{1.27} & \speedup{1.18} \\
    \midrule
    Σωρευτικά & \speedup{97.4} & \speedup{97.3} & \speedup{109.6} \\
    \bottomrule
  \end{tabular}
  \end{center}
  \vspace{0.3em}
  \footnotesize Πλατφόρμα: \en{BPI-F3} (\instret ablation), scalar \en{polymul} \HQC-128).
  % src: results/reports/ablation/polymul_scalar_zbc.md (instret incremental + cumulative)
\end{frame}

\begin{frame}{Παράρτημα: γιατί το \en{threshold} $T=16$ (και όχι $32/64$)}
  \begin{itemize}
    \item Το $T=16$ ταιριάζει το \VLMAX{} ($\VLEN{}=256$, \code{LMUL=m4}): ο leaf καλύπτει ακριβώς μία vector iteration
    \item Αύξηση $T$ αυξάνει τετραγωνικά τα γινόμενα ανά leaf ($T^2$ αντί $(T/2)^2\times3$), ακυρώνοντας ένα επίπεδο Karatsuba
  \end{itemize}
  \vspace{0.2em}
  \begin{center}
  \small
  \begin{tabular}{@{}lrr@{}}
    \toprule
    $T$ & median \instret & $\Delta$ vs $T=16$ \\
    \midrule
    \textbf{16} & \textbf{\num{354187}} & (βάση) \\
    32 & \num{365221} & \con{+\pct{3.1}} \\
    64 & \num{428584} & \con{+\pct{21.0}} \\
    128 & \num{540060} & \con{+\pct{52.5}} \\
    \bottomrule
  \end{tabular}
  \end{center}
  \vspace{0.2em}
\end{frame}

\begin{frame}{Παράρτημα: πλήρης πίνακας επιταχύνσεων \en{KEM} (\instret)}
  \small
  \begin{center}
  \begin{tabular}{lccc}
    \toprule
    & Tier 1 & Tier 2 & Tier 3 \\
    \midrule
    \HQC-128 & \speedup{35.0} & \speedup{55.9} & \speedup{75.9} \\
    \HQC-192 & \speedup{46.9} & \speedup{72.3} & \speedup{107.0} \\
    \HQC-256 & \speedup{60.1} & \speedup{86.9} & \speedup{163.8} \\
    \bottomrule
  \end{tabular}
  \end{center}
  \vspace{0.4em}
  \footnotesize \en{polymul} (\HQC-128): \speedup{101.3}/\speedup{124.6}/\speedup{264.1}.
  \HQC-256: \speedup{110.2}/\speedup{133.4}/\speedup{436.5}.
\end{frame}

\begin{frame}{Παράρτημα: \en{Keccak} \& αποκωδικοποίηση}
  \small
  \begin{center}
  \begin{tabular}{lccc}
    \toprule
    Επιτάχυνση & \HQC-128 & \HQC-192 & \HQC-256 \\
    \midrule
    \en{Keccak} & \speedup{1.36} & \speedup{1.36} & \speedup{1.36} \\
    \en{Decode} & \speedup{14.3} & \speedup{15.9} & \speedup{13.2} \\
    \bottomrule
  \end{tabular}
  \end{center}
\end{frame}

\begin{frame}{Παράρτημα: υλικό (BPI-F3, \HQC-128)}
  \small
  \begin{center}
  \begin{tabular}{lcc}
    \toprule
    & scalar (\Zbc) & vector (\RVV) \\
    \midrule
    \en{polymul} & \speedup{95.0} & \speedup{95.3} \\
    Πλήρες \en{KEM} & \speedup{37.5} & \speedup{45.9} \\
    \bottomrule
  \end{tabular}
  \end{center}
  \par\vspace{0.3em}
  \footnotesize \en{IPC}$\approx 1.49$, ο \instret υπερεκτιμά τους κύκλους κατά ${\approx}$\pct{50}.
\end{frame}

\begin{frame}{Παράρτημα: επικυρωμένο ISA vs custom hardware (kernel-level)}
  \centering
  \resizebox{\linewidth}{!}{% Kernel-level acceleration: ratified RISC-V ISA vs tightly-coupled custom HW.
% Present work: real BPI-F3 cycles (HQC-128) -> reduction & speedup vs reference kernel:
%   poly-mul  : hw-microbench-polymul-table (Reference 11635670 / Tier2 122085 = x95.31, 98.95%)
%   Keccak-f  : hw-microbench-keccak-table  (Reference 520235   / Tier1 356668 = x1.46,  31.4%)
% TYRCA (Dolmeta2025-TYRCA, Table III): Karatsuba 6072->52cc (99.14%, x116.8);
%   Keccak-f 26825->2538cc (90.54%, x10.57).
% Seo (Seo2026-HQC-Ibex, Table III): MULT_64 2444->31cc (98.73%, x78.84);
%   Keccak-f 25447->245cc (99.04%, x103.87).
\begin{thesistblr}{
  colspec = {l l r r},
  cell{2}{1} = {r=3}{l},
  cell{5}{1} = {r=3}{l},
  hline{5} = {0.4pt, figtablerule},
}
  {Πυρήνας} & {Προσέγγιση} & {Μείωση κύκλων} & {Speedup} \\
  Carry-less / poly-mul & Παρούσα εργασία (Tier~2) & \pct{98.95} & \speedup{95.31} \\
   & TYRCA (custom Karatsuba unit) & \pct{99.14} & \speedup{116.8} \\
   & Seo (custom \code{MULT\_64} ISE) & \pct{98.73} & \speedup{78.84} \\
  Keccak-f & Παρούσα εργασία (scalar \Zbb) & \pct{31.4} & \speedup{1.46} \\
   & TYRCA (custom Keccak unit) & \pct{90.54} & \speedup{10.57} \\
   & Seo (custom Keccak ISE) & \pct{99.04} & \speedup{103.87} \\
\end{thesistblr}
}
  % src: figures/ch9-discussion/ratified-vs-custom-hw-table.tex
  \pdfpcnote{Q alles RISC-V ergasies: TYRCA/Seo ftanoun x100+ se kernel ALLA me tightly-coupled custom ISE. Emeis: x95 poly-mul me MONO ratified Zbc/Zvbc, xoris custom silicon -> pio portable/maintainable.}
\end{frame}

\begin{frame}{Ευρύτερο \VLEN αποδίδει ελάχιστα: το κρίσιμο κλάσμα είναι οριακό}
  \HL{Μόνο $\approx\pct{3}$ του \KEM κλιμακώνει με ευρύτερο \VLEN.}
  \begin{columns}[c,onlytextwidth]
    \begin{column}{0.48\textwidth}
      \footnotesize
      \begin{itemize}
        \item Στο \VLEN$=\num{2048}$: Tier~3 $-\pct{2.6}$ (\HQC-256), Tier~2 $<\pct{1.7}$.
        \item Κυριαρχεί scalar εργασία (\en{Keccak}/\Zbb $\approx\pct{56.5}$).
        \item \code{vclmul\_leaf} κορεσμένο ($n=\num{16}$, $\VLMAX \ge n$).
        \item Μόνο Karatsuba (VLA strip-mining, $\approx\pct{2.9}$) κλιμακώνει.
      \end{itemize}
    \end{column}
    \begin{column}{0.52\textwidth}
      \centering
      \tikzsetnextfilename{pres-backup-vlen-budget-hbar}
      \begin{tikzpicture}
      \begin{axis}[
        width=\linewidth, height=5.0cm,
        xbar stacked, bar width=11pt,
        symbolic y coords={Other,Decode,Keccak,Polymul},
        ytick=data, xmin=0, xmax=60, xtick={0,20,40,60},
        xlabel={Ποσοστό \instret (\%)},
        legend style={at={(0.5,-0.30)}, anchor=north, legend columns=1,
                      font=\scriptsize, draw=none},
        every axis plot/.append style={draw=none},
        enlarge y limits=0.22,
      ]
      \addplot[fill=green!50!black,fill opacity=0.75] coordinates {(2.9,Polymul) (0,Keccak) (0.1,Decode) (0,Other)};
      \addlegendentry{Κλιμακώνει με \VLEN}
      \addplot[fill=blue!60,fill opacity=0.75] coordinates {(15.6,Polymul) (0,Keccak) (1.6,Decode) (0,Other)};
      \addlegendentry{Σταθερό (vectorized)}
      \addplot[fill=gray!50] coordinates {(10.1,Polymul) (56.5,Keccak) (5.3,Decode) (8.0,Other)};
      \addlegendentry{Scalar / μη κλιμακούμενο}
      \end{axis}
      \end{tikzpicture}
      % src: results/simulation/profiling/hotspots/vector_Zvk/hotspots_kem_L128_vector-zvk_latest.csv (commit bdd9cf4)
    \end{column}
  \end{columns}
  \pdfpcnote{HQC-128, Tier3, VLEN=256: scalable 3.0\% / fixed-vec 17.2\% / scalar 79.9\%. Horizontal layout -> kernel labels no overlap.}
\end{frame}
